\hypertarget{introduction}{%
\section{Introduction}\label{introduction}}

At the maximally mixed state, estimating purity from classical shadows is as expensive as it ever
gets. That sounds backwards. The estimand is \(2^{-n}\), nearly zero, and every snapshot returns
\(\mathrm{Tr}(G\rho) = 2^{-n}\) deterministically, so the first-order projection variance
\(\zeta_1\) vanishes identically. What does not vanish is the kernel term: the single-qubit second
moment \(\mathbb{E}[\mathrm{Tr}(G_1G_2)^2]\) equals \(7\) for a maximally mixed qubit, so
\(\zeta_2 = 7^n - 4^{-n}\), and the estimator's error is of order \(7^{n/2}/M\) against a target
shrinking as \(2^{-n}\). Low purity does not remove the exponential cost; it concentrates it.

The exponential is therefore a property of the estimator, and the question this paper answers is
what fixes its rate. The answer is not the purity, and not a constant of the protocol. It is two
spectral functionals of the state.

The finite-sample variance of the shadow U-statistic splits into projection terms of different
budget order --- the standard Hoeffding decomposition \cite{hoeffding1948, serfling1980}, stated
for shadows in \cite{huang2020shadows} --- and that the resulting decay rate changes with the
budget is established: Elben et al.~\cite{elben2020mixed} state both regimes and plot them, and
Aasen and G\"arttner \cite{aasen2026limitations} track the exponent migrating between them under
detector noise. What those analyses do not do is evaluate the coefficients. Bounding \(\zeta_1\)
and \(\zeta_2\) separately constrains their ratio not at all, and the exact second moment that
would fix it is applied in the literature to worst-case sample complexity rather than to a state.

\textbf{Thesis.} We evaluate the two projection variances exactly for an arbitrary state at
\(k = 2\), from the second-moment identity of \cite{huang2020shadows}: \(\zeta_2\) as a Pauli sum
with kernel \(14^{-|P|}\), \(\zeta_1\) as a cubic contraction of the same identity. Three things
follow. First, \(7^n - 1 \le \zeta_2 \le (17/2)^n\) for every state, so the per-qubit exponential
base is confined to \([7, 17/2]\) in the limit, and which value a state realizes is set by how its
Pauli weight is distributed rather than by its purity: a Haar-random pure state and a pure product
state, both of purity one, sit at \(7\) and at \(15/2\). Second, the budget at which the effective
error exponent crosses over is \(M^* = \zeta_2/(2\zeta_1)\), and its growth base factorizes as
\(\mathrm{base}(\zeta_2)/\mathrm{base}(\zeta_1)\) exactly, giving \(28/5\) for spread spectra and
exactly \(5\) for every pure product state. Third, because these are statewise functionals they
predict what an individual state does, which we verify by ordering states within a family whose
thresholds span a decade --- and, on a fixed-estimand control, show cannot be verified at all.

That the quantities are exact does not make them usable, since evaluating them needs the state.
We close that gap too, with unbiased estimators of both variances built from four disjoint blocks
of a pilot shadow budget, which never form the Pauli spectrum and whose cost grows more slowly in
\(n\) than the threshold they locate.

Set against two exact bias laws for collective measurement, the same machinery yields a plug-in
criterion, with nothing fitted to the crossover data, for the size at which a two-copy measurement
attains the lower RMSE at equal copy budget. That application is where the practical stakes sit,
and they are not small: on the noisy-pure states we benchmark, at a fixed budget and ten qubits,
the single-copy purity estimate projected into the interval purity is known a priori to occupy has
RMSE \(0.582\), while a constant fixed at the midpoint of that interval, which never looks at the
data, has \(0.310\). From seven qubits on, the single-copy route is worse than that constant.
Before projection the estimator is still unbiased; it is the spread that has become useless.

That matters because the quantity is not exotic. The moments \(\mathrm{Tr}(\rho^k)\) are the entry
point to much of what one wants to know about an unknown state: the second is the purity, their
logarithms give the R\'enyi entropies, moments of the partial transpose furnish PPT entanglement
criteria and bounds on negativity, and purity enters error-mitigation protocols such as virtual
distillation. Any experiment characterizing the state its device actually prepared needs these
numbers.

We compare two routes throughout. The single-copy route measures one copy at a time in randomized
local bases, building classical shadows \cite{huang2020shadows}, and forms an unbiased U-statistic
from the snapshots; it uses no entangling gates. The collective route holds \(k\) copies
simultaneously and measures the cyclic-permutation operator, whose expectation is exactly
\(\mathrm{Tr}(\rho^k)\); it has \(O(1)\) variance but requires entangling operations across
copies. The asymptotic separation between them is settled. Single-copy strategies require
exponentially more samples than collective ones for a family of learning tasks
\cite{huang2022advantage,chen2022separations}, and for protocols restricted to an identical
single-copy projection-valued measurement the lower bound for purity estimation is
\(\Omega(\max\{1/\varepsilon^2, 2^{n/2}/\varepsilon\})\) \cite{gong2024purity}. What is not settled
is the question an experimentalist faces: at this noise level and this system size, which route
has the lower RMSE at equal state-copy budget?

The gap is visible in recent work. A May 2026 study of quantum machine learning advantage on tens
of noisy qubits \cite{qmladvantage2026} declined to give its measure-first protocol multi-copy
access, citing both a result specific to its learning task and the resource demands of multi-copy
entangled measurements. That task is not moment estimation, but the caution about resource cost is
a premise we can test quantitatively. On the theory side, \cite{noisylearning2026} shows that the
exponential two-copy advantage for purity \emph{testing} collapses under local depolarizing noise;
their task is testing and ours is estimation, a distinction we take up in Section 8. The caution
is reasonable. What has not been available is a prediction of the size at which it stops applying.
